% !PRINTED 12
\weberpage{12}

If we set, as in § 2

\weq{5}{\begin{gathered}
a = (x_2x_3)(x_1x_4), \\
b = (x_3x_1)(x_2x_4), \\
c = (x_1x_2)(x_3x_4), \\
a + b + c = 0,
\end{gathered}}

then, since only the ratio of $\alpha$ to $\beta$ matters, we can satisfy the last of the equations (4) by setting: % [sic] German source prints "könnnen" with three n-strokes

\weq{6}{\sqrt{\alpha} = \sqrt{\pm a} + \sqrt{\mp b}, \quad \sqrt{\beta} = \sqrt{\pm a} - \sqrt{\mp b},}

and from (3) there results for $z$ the expression

\weq{7}{z = \sqrt{\alpha}\,\dfrac{(xx_2)\sqrt{(x_3x_1)(x_4x_1)} + (xx_1)\sqrt{(x_3x_2)(x_4x_2)}}{(xx_2)\sqrt{(x_3x_1)(x_4x_1)} - (xx_1)\sqrt{(x_3x_2)(x_4x_2)}}.}

If beside (3) one sets up also the equation which follows from it:

\[
h'\,\dfrac{x-x_3}{x-x_4} = \dfrac{z-\sqrt{\beta}}{z+\sqrt{\beta}}
\]

then by logarithmic differentiation there results:

\[
\dfrac{(x_1x_2)\,dx}{(xx_1)(xx_2)} = \dfrac{2\sqrt{\alpha}\,dz}{z^2-\alpha}, \quad \dfrac{(x_3x_4)\,dx}{(xx_3)(xx_4)} = \dfrac{2\sqrt{\beta}\,dz}{z^2-\beta},
\]

and from this by multiplication, with regard to (5) and (6):

\weq{8}{\dfrac{dx}{\sqrt{\pm(xx_1)(xx_2)(xx_3)(xx_4)}} = \dfrac{2\,dz}{\sqrt{(z^2-\alpha)(z^2-\beta)}}.}

If now the four roots of $f(x)$ are real, then, as we have seen in § 2, there are eight ways of assigning these roots to the symbols $x_1$, $x_2$, $x_3$, $x_4$ in such a manner that $a$ and $b$ have opposite signs; and if therefore $\pm a$, $\mp b$ are positive, then $\sqrt{\alpha}$, $\sqrt{\beta}$ become real, hence $\alpha$, $\beta$ positive, and by (7) $z$ then becomes real as well.

If, secondly, $x_1$, $x_2$ are real and $x_3$, $x_4$ conjugate imaginary, then $a$ and $-b$ are conjugate imaginary; hence, if the signs of the square roots $\sqrt{\pm a}$, $\sqrt{\mp b}$ are suitably determined, $\sqrt{\alpha}$ becomes real and $\sqrt{\beta}$ purely imaginary, hence $\alpha$ positive, $\beta$ negative and $z$ real [because $(x_3x_1)(x_4x_1)$ and $(x_3x_2)(x_4x_2)$, as products of conjugate imaginary quantities, are positive].

If, finally, $x_1$, $x_2$ and $x_3$, $x_4$ are two conjugate imaginary pairs, then $a$ and $-b$ are both positive. If we therefore take the lower signs in (6), then $\sqrt{\alpha}$, $\sqrt{\beta}$ become purely imaginary, hence $\alpha$ and $\beta$ negative, and from (7) a real expression results for $z$.
